\scanpage{92}
\mspage{85}
% archive sheet number: 91
\[
\operatorname{card}\widetilde I\leq(\unsure{\operatorname{Card}|\operatorname{card}}I)^{\pi}\leq\pi'^\pi\qquad\text{si }\operatorname{card}I\leq\pi'^\pi
\]
(notamment, si $X\in\mathrm{Filt}^{\pi'}$ et qu'on prenne $I$ comme indiqué dans la condition c des filtr. cardinales). De plus, cette fois, le fait que $\widetilde I$ soit grand devant $\pi$ est trivial, on a même mieux, \inserted{\unsure{: pourvu que} $I$ soit filtrant:} car \uline{Toute} partie \inserted{$(I'_\alpha)_{\alpha\in A}$} de $\widetilde I$ \inserted{filtrante et} de cardinal $\leq\pi$ admet une borne supérieure \deleted{\unsure{$I$}} $\tilde i_A$ (qui est bien filtrante).

savoir la \inserted{réunion} \deleted{\unsure{partie engendrée et|partie engendrée|partie} filtrante de $I$ engendrée par \inserted{\unsure{on prend} $I'_0$} $I'_0=\bigcup_{\alpha\in A}I'_\alpha$, \unsure{qui est une} partie de cardinal $\leq\pi^2=\pi$, et qui \unsure{la|est une} partie filtrante}

De plus, le système inductif $(\unsure{X|\alpha}_{\tilde i})_{\tilde i\in\widetilde I}$ a la propriété plus forte que
\[
X_{\tilde i_A}\xleftarrow{\sim}\varprojlim_{\tilde i\in A}X_{\tilde i}.
\]
Ainsi on a la

\uline{Proposition 1 bis} Supposons \deleted{$\mathrm{Filt}^{\pi}$ stable} \inserted{$E$ stable par petites} $\unsure{\varinjlim|\underline{\lim}}$ filtrantes, et $\mathrm{Filt}^{\pi}$ stable par celles-ci. \deleted{Alors} si $\pi\geq\pi_0$ est donné. Alors pour tt $X$ dans $E$, on peut trouver un syst. ind. $(X_i)_{i\in I}$ de $\varinjlim_I X$, avec

a) $X_i\in\mathrm{Filt}^{\pi}$ pour tt $i\in I$, $I$ ens. ordonné grd dev. à $\pi$

b) Toute partie $J$ de $I$ filtrante de card. $\leq\pi$ admet une borne sup. $i_J$ dans $I$,

c) \unsure{dans c)|dans ce}, sous les conditions
